\begin{table}[!htbp]
\centering
\footnotesize
\caption{Support and selection diagnostics for the film-visibility model.}
\label{tab:si-visibility-selection}
\renewcommand{\arraystretch}{1.05}
\setlength{\tabcolsep}{3pt}
\begin{tabular}{>{\raggedright\arraybackslash}p{0.21\linewidth}>{\centering\arraybackslash}p{0.06\linewidth}>{\centering\arraybackslash}p{0.055\linewidth}>{\centering\arraybackslash}p{0.05\linewidth}>{\centering\arraybackslash}p{0.055\linewidth}>{\centering\arraybackslash}p{0.05\linewidth}>{\centering\arraybackslash}p{0.085\linewidth}>{\centering\arraybackslash}p{0.10\linewidth}}
\toprule
Time-respecting film-linked observations & N obs. & N songs & \makecell[b]{Mean\\ age} & \makecell[b]{Mean\\ weeks} & \makecell[b]{Mean\\ peak} & \makecell[b]{Mean\\ realized\\ films} & \makecell[b]{Mean\\ popularity} \\
\midrule
without recorded worldwide box office & 1,301 & 430 & 44.3 & 12.9 & 33.9 & 1.06 & 40.9 \\
with recorded worldwide box office & 15,245 & 4,294 & 38.7 & 15.9 & 26.5 & 2.55 & 50.3 \\
\bottomrule
\end{tabular}
\begin{flushleft}\footnotesize
The universe is every song-snapshot linked to at least one film released by the snapshot year. Rows split this universe by whether positive worldwide box office is recorded for at least one eligible film. Recorded totals may include receipts after the snapshot date. The row with observed box office is exactly the estimation sample of the continuous visibility model and its tercile visualization (15,245 song-snapshots, 4,294 normalized song-artist identities); the model's complete-case condition drops no observations.
\end{flushleft}
\end{table}
